\documentclass[aps,pre,reprint,nofootinbib]{revtex4-2}

\usepackage{amsmath,amssymb,amsthm,bm}
\usepackage{array}
\usepackage{booktabs}
\newcolumntype{L}[1]{>{\raggedright\arraybackslash}p{#1}}
\usepackage{graphicx}
\usepackage{hyperref}
\usepackage{microtype}
\usepackage{pgfplots}
\usepgfplotslibrary{groupplots}
\pgfplotsset{compat=1.15,
  paperaxis/.style={
    width=\columnwidth, height=0.58\columnwidth,
    grid=major, grid style={gray!20},
    tick label style={font=\footnotesize},
    label style={font=\footnotesize},
    legend style={font=\footnotesize, draw=none, fill=none},
    line width=0.9pt}}

\newtheorem{theorem}{Theorem}
\newtheorem{definition}{Definition}
\newcommand{\M}{\mathcal{M}}
\newcommand{\X}{\mathcal{X}}
\newcommand{\dd}{\mathrm{d}}

\begin{document}

\title{Entropy Across Singular Projections:\\
Fiber Multiplicity, Caustics, and Observer-Relative Information}

\author{Andrew H. Bond}
\email{andrew.bond@sjsu.edu}
\thanks{ORCID: 0009-0003-2599-6158}
\affiliation{San Jos\'e State University, San Jos\'e, California, USA}
\date{August 4, 2026}

\begin{abstract}
A lower-dimensional observer who watches a hidden dynamical system through a
projection must count, distinguish, and assign probability to the hidden
states compatible with each observation.  When the map from the hidden
evolution parameter to observed time has singularities, that counting changes
abruptly.  This paper treats the simplest such singularity, the fold, at
finite observational resolution.  We prove a local theorem stating that a
nondegenerate fold changes the number of hidden preimages by two while an
orientation-weighted count is conserved.  We compute closed-form branch
entropies for a family of canonical control maps.  The quadratic fold carries
exactly one bit at every two-branch slice, the symmetric slice of a double
fold carries exactly one and one half bits, and a degenerate critical point
of a monotone map carries no entropy at all, so multiplicity rather than
criticality is what stores the information.  We compute the resolution
scaling of the binned entropy of the fold caustic and verify numerically that
it converges to its differential limit with deviations shrinking as the
square root of the bin width.  Finally we demonstrate, in a reversible
Hamiltonian model, that the observed binned entropy of the coordinate-time
marginal changes by more than two bits along a trajectory while momentum
reversal recovers the initial hidden ensemble to near machine precision, so
the observed entropy change is produced and destroyed by folds of the
projection with zero hidden information loss.  Nothing in this paper is a
claim about physical spacetime, about the thermodynamics of real systems, or
about quantum gravity.
\end{abstract}

\maketitle

\section{Introduction}
\label{sec:intro}

An observer rarely sees the full state of the system it observes.  It sees a
projection, and it parameterizes what it sees by its own time.  If the hidden
evolution parameter and the observed time increase together everywhere, the
observer's bookkeeping is simple.  Each observed instant corresponds to one
hidden state, and observed uncertainty tracks hidden uncertainty.  The
interesting case is the other one.  If the observed time, viewed as a
function of the hidden evolution parameter, has a critical point, then a
single observed instant can correspond to several hidden states at once.
The set of hidden states compatible with one observation, which we call the
fiber of that observation, changes its size discontinuously as the observed
time crosses the critical value.

This paper studies how a finite-resolution observer counts and weighs the
members of such fibers, and how much information the observer's ignorance is
worth, in the simplest singular situation.  That situation is the fold, the
generic critical point of a map from one dimension to one dimension.
Singularity theory guarantees that folds are stable under small
perturbations, so they are the correct first object of study, and it
classifies everything more complicated as a higher catastrophe.

Four results are reported, and each carries an explicit evidence label.
First, a local theorem, proved, describing what a fold does to the fiber.
The number of hidden preimages of an observed instant jumps by two across a
fold, while a signed count that weights each preimage by its time
orientation does not change at all.  Second, closed-form entropy values,
computed and verified in two independent implementations, for a family of
canonical control maps.  A quadratic fold is worth exactly one bit of branch
ambiguity, a symmetric double fold is worth exactly one and one half bits,
and a monotone map with a degenerate critical point is worth nothing.
Third, a resolution-scaling law, computed and numerically demonstrated, for
the observed density near a fold.  That density diverges at the fold, but
the divergence is integrable, and the binned observed entropy converges to a
finite differential limit as the bin width shrinks, with deviations falling
as the square root of the bin width.  Fourth, a numerical demonstration in a
small Hamiltonian model.  A reversible hidden dynamics, integrated with a
time-symmetric method, produces observed entropy changes of more than two
bits along a trajectory, and then gives every bit back.  Reversing the
momenta recovers the initial hidden ensemble to near machine precision, and
the entropy curve retraces itself exactly.  The observed entropy change is
therefore a property of the projection and the observer's resolution, not of
any loss in the hidden dynamics.

The scope of these claims is deliberately narrow.  The theorem is a
statement in one-dimensional singularity theory.  The closed-form values are
exact computations for named control maps.  The demonstration is a
simulation of a toy model.  No claim is made about physical spacetime, about
the thermodynamics of any real system, or about quantum gravity.  The paper
does establish, within its model class, a clean separation between two
things that are often conflated, the growth of an observer's coarse-grained
entropy and the destruction of hidden information.  Folds of a projection
produce the first without any of the second.

\section{Setup}
\label{sec:setup}

Let $\M$ be a smooth manifold, the hidden state space, and let
\begin{equation}
  z: I \rightarrow \M
\end{equation}
be a smooth trajectory parameterized by a hidden evolution variable $\tau$
on an interval $I$.  A smooth observation map
\begin{equation}
  \pi: \M \rightarrow \X
\end{equation}
sends hidden states to an observed space $\X$, which we take to carry
spacetime coordinates $x^\mu$.  An observer supplies a time functional
\begin{equation}
  T_u(x) = u_\mu x^\mu,
\end{equation}
with $u$ the observer's time direction, and reads the trajectory through
the scalar time map
\begin{equation}
  f(\tau) = T_u\!\left(\pi(z(\tau))\right).
  \label{eq:timemap}
\end{equation}

\begin{definition}[Fiber of an observed time]
The fiber of the observed time $t$ is the preimage
\begin{equation}
  F(t) = \{\, \tau \in I : f(\tau) = t \,\},
\end{equation}
the set of hidden parameter values compatible with the observation $t$.
\end{definition}

The unsigned multiplicity $N(t) = |F(t)|$ counts the fiber.  The signed
count weights each element by its time orientation,
\begin{equation}
  N_{\mathrm{sgn}}(t)
  = \sum_{\tau_i \in F(t)} \operatorname{sgn} f'(\tau_i),
  \label{eq:signed}
\end{equation}
defined on regular slices, meaning values of $t$ whose fiber contains no
critical point of $f$.

The observer does not resolve $t$ exactly.  It resolves $t$ to within a bin
width $\epsilon > 0$, and its observed distribution is the pushforward of
the hidden distribution under $f$, accumulated into bins of width
$\epsilon$.  We write $H_\epsilon$ for the Shannon entropy in bits of that
binned distribution.  Finite resolution is mandatory rather than a modeling
convenience.  Differential entropy is not invariant under reparameterization
of the observed variable, and near a singularity of $f$ the pushforward
density diverges, so an unbinned continuous entropy assigned to the observed
variable depends on coordinate choices in exactly the region this paper is
about.  The binned quantity $H_\epsilon$ is well defined for every
$\epsilon > 0$, and every entropy statement below is either a statement
about $H_\epsilon$ at stated $\epsilon$, a statement about a discrete
branch-label distribution, or an exact limit as $\epsilon \to 0$ with the
divergent resolution term removed.

\section{The fold theorem}
\label{sec:fold}

The fold is the nondegenerate critical point of the time map.  The
statement below is kinematic.  It concerns the map $f$ alone and involves
no dynamics, no probability, and no physics.

\begin{theorem}[Local fold multiplicity]
\label{thm:fold}
Let $f$ be smooth near $\tau_0$ with $f'(\tau_0) = 0$ and
$f''(\tau_0) \neq 0$, and let $t_0 = f(\tau_0)$.  Then in suitable smooth
local coordinates
\begin{equation}
  f(\tau) - t_0 = \sigma\,(\tau - \tau_0)^2,
  \qquad \sigma \in \{+1, -1\}.
  \label{eq:normalform}
\end{equation}
For $\sigma = +1$ the unsigned multiplicity of nearby regular slices is $0$
for $t < t_0$ and $2$ for $t > t_0$, and the two branches carry opposite
signs of $f'$.  The signed count $N_{\mathrm{sgn}}$ takes the same value on
both sides of the fold.  For $\sigma = -1$ the time-reversed statement
holds.
\end{theorem}

\begin{proof}
The one-dimensional Morse lemma gives the normal form
\eqref{eq:normalform}.  For $\sigma = +1$ and $t > t_0$ the local preimages
are $\tau_\pm = \tau_0 \pm \sqrt{t - t_0}$, with
$f'(\tau_\pm) = \pm 2\sqrt{t - t_0}$, so the two branches have opposite
orientation and their contributions to \eqref{eq:signed} cancel.  For
$t < t_0$ there are no local preimages, contributing zero.  The local signed
count is therefore $0$ on both sides.  The case $\sigma = -1$ follows by
reversing the sign of the time coordinate.
\end{proof}

The fold is stable.  Whitney's classification of stable maps of the plane
and the catastrophe-theory literature descending from it establish that a
nondegenerate fold persists under small smooth perturbations and that its
square-root branch geometry is universal \cite{Whitney1955,Arnold1992}.
Those structural statements are standard mathematics and are cited, not
reproved.  What the present campaign adds is a numerical verification of
the same statements by the instrument stack used throughout this paper, and
that verification is labeled demonstrated, not proved.  In five hundred
perturbation trials per amplitude, the fold persisted with probability
$1.000$ at every perturbation norm up to $0.12$.  The branch-separation
exponent, fitted over six decades of distance to the fold, was
$0.500000000000$ against the square-root law, with fit error below
$10^{-10}$.  The deviations of the perturbed fold location and fold
curvature from first-order perturbation theory scaled as the square of the
perturbation amplitude across seven decades of measured deviation.  At the
largest tested amplitude, norm $0.32$, persistence dropped to $0.956$,
which is the expected exit from the perturbative regime rather than an
instrument failure.  These numbers validate the instruments against known
mathematics.  They are not new mathematics.

\section{Branch weights and branch entropy}
\label{sec:weights}

Suppose the hidden parameter $\tau$ is distributed uniformly on $I$.  The
coarea formula for the pushforward of a density under a map with
nonvanishing derivative on the fiber assigns each preimage
$\tau_i \in F(t)$ a weight proportional to the reciprocal of the local time
speed \cite{EvansGariepy1992},
\begin{equation}
  w_i(t) = \frac{1/|f'(\tau_i)|}{\sum_{j} 1/|f'(\tau_j)|}.
  \label{eq:coarea}
\end{equation}
An observer who resolves the observed time but not the branch label is
ignorant of which preimage the system occupies, and the value of that
missing label is the branch entropy
\begin{equation}
  H_{\mathrm{br}}(t) = -\sum_i w_i(t)\, \log_2 w_i(t)
  \label{eq:branchentropy}
\end{equation}
in bits \cite{CoverThomas2006,Jaynes1957}.

The following values are proved by direct computation from
\eqref{eq:coarea} and \eqref{eq:branchentropy} and were verified in two
independent implementations, one in MATLAB and one in Python, which agree
on branch locations to $4.4 \times 10^{-16}$.

\emph{Quadratic fold.}  For $f(\tau) = \tau^2$ the two branches at any
regular slice $t > 0$ sit at $\tau = \pm\sqrt{t}$ with equal speeds
$|f'| = 2\sqrt{t}$.  The weights are $(\tfrac12, \tfrac12)$ and
\begin{equation}
  H_{\mathrm{br}}(t) = 1 \text{ bit}
  \qquad \text{for every } t > 0.
\end{equation}
The fold hands the observer exactly one new binary distinction, uniformly
across its two-branch side.

\emph{Double fold.}  For $f(\tau) = \tau^3 - \tau$, at the symmetric slice
$t = 0$ the fiber is $\{-1, 0, 1\}$ with speeds $|f'| = (2, 1, 2)$.  The
weights are $(\tfrac14, \tfrac12, \tfrac14)$ and
\begin{equation}
  H_{\mathrm{br}}(0) = \tfrac{3}{2} \text{ bits}
\end{equation}
exactly.  As the slice approaches a band edge, the two merging branches
dominate the fiber measure because their speeds vanish, and the branch
entropy tends to $1$ bit.  The measured deviation from $1$ bit is below
$10^{-3}$ at a distance of $10^{-9}$ from the edge.

\emph{Monotone maps.}  For $f(\tau) = \tau$, and equally for the degenerate
cubic $f(\tau) = \tau^3$, every fiber is a single point and
\begin{equation}
  H_{\mathrm{br}}(t) = 0 \text{ bits}
  \qquad \text{for every regular } t.
\end{equation}
The cubic has a critical point at the origin, yet contributes no branch
entropy anywhere, because the map remains monotone and no slice ever meets
more than one preimage.  Multiplicity, not criticality, is what carries the
information.  Throughout all four control maps the signed count
\eqref{eq:signed} is constant on regular slices, so the observational
multiplicity and its entropy grow while the topological invariant is
conserved, exactly as Theorem~\ref{thm:fold} requires.

\begin{figure}[t]
\centering
\begin{tikzpicture}
\begin{loglogaxis}[paperaxis, height=0.5\columnwidth,
  xlabel={bin width $\varepsilon$},
  ylabel={deviation from limit (bits)},
  legend pos=north west]
\addplot[only marks, mark=*, mark size=2pt, blue!70!black]
  table[x=eps, y=deviation] {figdata/pe1_convergence.dat};
\addlegendentry{measured}
\addplot[dashed, domain=1e-6:1e-4, samples=2, black!60]{0.4472*sqrt(x)};
\addlegendentry{$\propto\sqrt{\varepsilon}$}
\end{loglogaxis}
\end{tikzpicture}
\caption{Convergence of the binned fold-caustic entropy to its
differential limit.  The deviation of $H_{\varepsilon}+\log_{2}\varepsilon$
from $1-1/\ln 2$ bits shrinks proportionally to the square root of the bin
width across the tested range, so the inverse-square-root caustic is
integrable and produces no entropy pathology.  Points are recomputed
deterministically from the sealed instrument.}
\label{fig:convergence}
\end{figure}

\section{The fold caustic and resolution scaling}
\label{sec:caustic}

Where the derivative of the time map vanishes, the projection compresses
hidden measure into a small observed interval, and the observed density
diverges.  For the quadratic fold $f(\tau) = \tau^2$ with $\tau$ uniform on
$[-1, 1]$, the pushforward density is
\begin{equation}
  p(t) = \frac{1}{2\sqrt{t}}, \qquad 0 < t \le 1,
  \label{eq:caustic}
\end{equation}
the inverse-square-root divergence familiar from caustics in optics, where
a family of rays folds onto an envelope and the intensity diverges with the
same exponent at the same codimension \cite{BerryUpstill1980}.  The
divergence in \eqref{eq:caustic} is integrable, so the fold concentrates
observed probability without creating any entropy pathology.

The binned observed entropy of \eqref{eq:caustic} obeys a clean resolution
law.  A direct computation of the differential entropy of \eqref{eq:caustic}
gives $1 - 1/\ln 2$ bits, approximately $-0.4427$ bits, the negative value
being ordinary for a density concentrated on a subunit interval.  The
binned entropy at width $\epsilon$ then satisfies
\begin{equation}
  H_\epsilon + \log_2 \epsilon \;\longrightarrow\; 1 - \frac{1}{\ln 2}
  \qquad (\epsilon \to 0),
  \label{eq:scaling}
\end{equation}
with the bin masses computed exactly from the cumulative distribution
$\sqrt{t}$.  The measured deviations from the limit are $0.00447$ at
$\epsilon = 10^{-4}$, $0.00141$ at $\epsilon = 10^{-5}$, and $0.000447$ at
$\epsilon = 10^{-6}$, shrinking as $\sqrt{\epsilon}$, the rate set by the
single bin that straddles the divergence.  The convergence
\eqref{eq:scaling} is the quantitative sense in which the caustic is
harmless.  All resolution dependence sits in the universal term
$-\log_2 \epsilon$ shared by every absolutely continuous observed density,
and the fold contributes only a finite, computable offset.  This
convergence was verified in the stated range and is labeled demonstrated
there, with the closed-form limit computed exactly.

\begin{figure}[t]
\centering
\begin{tikzpicture}
\begin{groupplot}[group style={group size=1 by 2, vertical sep=5mm,
    xlabels at=edge bottom, xticklabels at=edge bottom},
  paperaxis, height=0.42\columnwidth,
  xlabel={observation level $t$}]
\nextgroupplot[ylabel={preimage count}, ymin=3, ymax=14, ytick={4,8,12}]
\addplot[blue!70!black, const plot]
  table[x=level, y=multiplicity] {figdata/pe2_staircase.dat};
\nextgroupplot[ylabel={entropy (bits)}, ymin=0]
\addplot[red!70!black]
  table[x=level, y=entropy] {figdata/pe2_staircase.dat};
\end{groupplot}
\end{tikzpicture}
\caption{The observation axis of one hidden trajectory of the toy
Hamiltonian.  The unsigned preimage count (top) is a staircase whose
breakpoints are the fold values of the observed time and whose interior
band profile is $4$, $8$, $12$, $13$ before descending symmetrically.  The
coarea branch entropy (bottom) rises and falls with the staircase.  The
signed count, not shown, equals the path degree at every tested level.}
\label{fig:staircase}
\end{figure}

\begin{figure}[t]
\centering
\begin{tikzpicture}
\begin{axis}[paperaxis, height=0.5\columnwidth,
  xlabel={hidden evolution parameter $\tau$},
  ylabel={observed entropy $H_{\varepsilon}$ (bits)},
  ymin=0, legend pos=south east]
\addplot[blue!70!black]
  table[x=tau, y=forward] {figdata/pe2_cycle.dat};
\addlegendentry{forward}
\addplot[only marks, mark=o, mark size=1.5pt, each nth point=8,
  red!70!black]
  table[x=tau, y=reverse_reflected] {figdata/pe2_cycle.dat};
\addlegendentry{reversed, reflected}
\end{axis}
\end{tikzpicture}
\caption{Observed binned entropy of the coordinate-time marginal for the
ten-thousand-member ensemble.  The forward curve spans $2.07$ bits.  The
open circles show the momentum-reversed run reflected in time, which
retraces the forward curve with defect exactly zero while the initial
hidden ensemble is recovered to $1.1\times 10^{-14}$, so the entire
entropy change is observational.}
\label{fig:cycle}
\end{figure}

\section{A reversible fold cycle}
\label{sec:cycle}

The results above are static.  This section demonstrates, in a model, what
folds do to an observer's entropy along a dynamical trajectory, and what
they do not do.  Everything in this section carries the label demonstrated
in a model, at exploratory grade, and nothing in it is a thermodynamic
claim.

The hidden dynamics is the conserved toy Hamiltonian
\begin{equation}
\begin{split}
  H = {}& \tfrac12\!\left(p_t^2 + p_x^2 + p_u^2\right)
  + \tfrac12 \omega_t^2 t^2
  + \tfrac12 \omega_u^2 u^2 \\
  & + \tfrac{\lambda}{4} u^4
  + g E\, t u,
\end{split}
\label{eq:toy}
\end{equation}
with $\omega_t = 1$, $\omega_u = 1.2$, $\lambda = 0.1$, $g = 0.25$, and
$E = 1$.  The coordinate $t$ plays the role of observed time, with
$\dd t/\dd\tau = p_t$, the coordinate $u$ is hidden, and folds of the time
map occur wherever $p_t$ crosses zero with nonzero acceleration.  The
integrator is velocity Verlet with step $10^{-3}$ over
$\tau \in [0, 40]$, a symplectic and time-symmetric method, and the
relative energy drift over the full run was $6.1 \times 10^{-9}$, within the sealed
fixed-step bar of the tolerance freeze described in
Section~\ref{sec:methods}.  The model \eqref{eq:toy} is not relativistic
and is not proposed as a theory of anything.  It is an instrumented
example whose only job is to produce genuine folds under a genuinely
reversible dynamics.

\emph{Single-trajectory audit.}  One trajectory of \eqref{eq:toy} produces
twelve folds of the time map.  Observed through constant-time levels, its
unsigned multiplicity is a staircase over the observation axis with band
profile $4, 8, 12, 13$ and then the same values descending, constant
between breakpoints, and the signed count \eqref{eq:signed} obeys the path
degree rule at each of the roughly $370$ tested levels.  That rule states
that the signed crossing count of a level equals the sign of the net
displacement of $t$ when the level lies between the endpoint values of $t$
and equals zero otherwise.  One caveat is recorded with this result rather
than elsewhere.  The twelve fold values of $t$ cluster near the turning
amplitudes of the trajectory, so only the two isolated breakpoints, the
endpoint values with steps of one, admitted the per-breakpoint step check
at this level resolution, while the clustered fold breakpoints were
verified in aggregate through steps of four across close pairs together
with band constancy.  A finer level grid around the clusters is the
designed follow-up.

\emph{Reversible ensemble.}  An ensemble of $10^4$ members, initialized
from a fixed seeded Gaussian spread and integrated with the same method,
gives the observed side of the demonstration.  The observed quantity is the
binned entropy $H_\epsilon$ of the coordinate-time marginal, with bin width
$0.05$ on the range $[-3, 3]$.  Along the forward trajectory that entropy
spans $2.0708$ bits, from $0.1952$ bits to $2.2660$ bits, rising and
falling as fold bands raise and lower the fiber multiplicity.  The hidden
side of the demonstration is the reversal.  Flipping every momentum at the
final time and integrating forward again recovers the initial ensemble
with maximum residual $1.1 \times 10^{-14}$ over all members and
coordinates, and the reversed entropy curve retraces the forward curve
with defect exactly zero, a bit-identical retrace guaranteed by the
time symmetry of the integrator.  The falsification bar for this run,
sealed in advance, would have been tripped by any signed-count violation,
any failure of band constancy, a recovery residual above $10^{-9}$, or a
reversed entropy curve that failed to retrace the forward one.  None
tripped.

The conclusion is stated exactly.  In this model, more than two bits of
observed entropy appear and disappear along a trajectory whose hidden
dynamics loses nothing.  The observed entropy change is observational.  It
is produced and destroyed by folds of the projection together with the
observer's finite resolution, with zero hidden information loss.

Three distinct quantities must not be conflated in reading that sentence.
The first is the conditional entropy of the hidden state given the
observation, the observer's residual uncertainty, which folds change by
changing the number and weighting of compatible hidden states.  The second
is the coarse-grained entropy of the observed, binned distribution, the
quantity measured above, which a reversible hidden dynamics can raise and
lower freely.  The third is thermodynamic entropy production, which
requires a physical mechanism that makes lost distinctions inaccessible,
such as dissipation, an unmonitored environment, or the erasure of a
memory, and which carries the Landauer cost
\cite{Landauer1961,Bennett1982}.  Nothing in this paper measures or bounds
the third quantity.  The demonstration shows precisely that the first two
can undergo large changes while the third has no mechanism to occur, since
the dynamics is reversible and reversed.

\section{Entropy signatures of projection singularities}
\label{sec:signatures}

The results above organize into a classification of what each local type of
projection singularity contributes to an observer's information budget.
Table~\ref{tab:signatures} states it.  The fold row and the two regular
rows are established by the preceding sections at the labels given there.
The cusp and higher rows are design, not results.  They state the
predictions this framework makes and the experiments that would test them,
and no measurement reported here bears on them.

\begin{table}[t]
\caption{\label{tab:signatures}Entropy signatures by singularity class.
Rows above the rule are established in this paper at the stated labels.
Rows below the rule are open questions stated as design targets.}
\begin{ruledtabular}
\begin{tabular}{L{0.14\columnwidth}L{0.2\columnwidth}L{0.15\columnwidth}L{0.24\columnwidth}}
Class & Multiplicity & Signed count & Entropy signature \\
\colrule
Regular & constant & conserved & none \\
Degenerate monotone & constant & conserved & none \\
Fold & jumps by two & conserved & one new binary distinction \\
\colrule
Cusp & varies across fold boundaries & conserved globally & correlated branch ambiguity (open) \\
Higher catastrophes & larger structured changes & constrained & multiway ambiguity (open) \\
\end{tabular}
\end{ruledtabular}
\end{table}

The organizing statement is the pairing of the fold row with the
degenerate-monotone row.  A fold gives the observer exactly one new binary
distinction, worth one bit at equal weighting and never more than one bit,
while a degenerate critical point of a monotone map gives the observer
nothing at all.  Criticality of the time map is therefore neither
necessary nor sufficient for observational entropy.  What is both
necessary and sufficient, in the one-dimensional setting treated here, is
a change in fiber multiplicity, and Theorem~\ref{thm:fold} fixes the
quantum of that change at two while conserving the signed count.  Whether
the cusp contributes a correlated two-bit structure, and whether the
higher catastrophes contribute the multiway ambiguities their normal
forms suggest, are the designed continuations, and they are listed in the
table as open.

\section{Methods and reproducibility}
\label{sec:methods}

All instruments used in this paper were validated against exact analytic
controls before any result-bearing run, under a sealed tolerance freeze
registered as PF0-FREEZE-001, which fixes every numerical bar in advance
and voids itself if edited.  The control net comprises a monotone null, the
exact quadratic creation and annihilation folds, the degenerate cubic, and
the double fold, in both a MATLAB implementation and an independent Python
implementation.  Cross-language branch locations agree to
$4.4 \times 10^{-16}$, and the same scenario executed on two substrates
yields bit-identical scenario and outcome hashes.  Every result-bearing run
writes an append-only JSON evidence record in the repository containing the
complete parameter set, seeds, outcome arrays, and SHA-256 hashes of inputs
and outputs, and the records cited here carry the exploratory label of the
campaign's evidence-grading scheme.  A reducer may summarize records but
may not rewrite them.

\begin{thebibliography}{9}

\bibitem{Whitney1955}
H.~Whitney,
On singularities of mappings of Euclidean spaces. I. Mappings of the plane
into the plane,
Ann.\ of Math.\ \textbf{62}, 374 (1955).

\bibitem{Arnold1992}
V.~I.~Arnold,
\emph{Catastrophe Theory}, 3rd ed.\
(Springer-Verlag, Berlin, 1992).

\bibitem{BerryUpstill1980}
M.~V.~Berry and C.~Upstill,
Catastrophe optics: morphologies of caustics and their diffraction patterns,
Prog.\ Opt.\ \textbf{18}, 257 (1980).

\bibitem{Jaynes1957}
E.~T.~Jaynes,
Information theory and statistical mechanics,
Phys.\ Rev.\ \textbf{106}, 620 (1957).

\bibitem{CoverThomas2006}
T.~M.~Cover and J.~A.~Thomas,
\emph{Elements of Information Theory}, 2nd ed.\
(Wiley-Interscience, Hoboken, 2006).

\bibitem{EvansGariepy1992}
L.~C.~Evans and R.~F.~Gariepy,
\emph{Measure Theory and Fine Properties of Functions}
(CRC Press, Boca Raton, 1992).

\bibitem{Landauer1961}
R.~Landauer,
Irreversibility and heat generation in the computing process,
IBM J.\ Res.\ Dev.\ \textbf{5}, 183 (1961).

\bibitem{Bennett1982}
C.~H.~Bennett,
The thermodynamics of computation, a review,
Int.\ J.\ Theor.\ Phys.\ \textbf{21}, 905 (1982).

\end{thebibliography}

\end{document}
